%!TEX root = ../main.tex
%
% PARKED 2026-09-20. This was section 3 of the Class 5 ISAC lecture: the ground-up
% OFDM primer. It was on the slides for the 9/9 delivery but never actually taught --
% that lecture used OFDM as a radar and skipped the machinery entirely.
%
% The frames themselves are NOT lost: all seventeen were moved into
% slides/Class_8_WiFi_History/03_first_speed_race.tex, where they now run before the
% 802.11a slide and are the course's only OFDM instruction. What is kept here is the
% ISAC-specific framing around them -- the "Question 2: which signal do we actually
% use?" roadmap frame, which only makes sense next to the three-questions frame in
% 02_dfrc.tex. The one frame that stayed in the ISAC deck, "OFDM already carries a
% radar inside it", now opens 05_applications.tex.
%
% Not input by main.tex.

%!TEX root = ./main.tex

\section[OFDM]{The Waveform We Already Have}

% SECTION TIMING: 32:00--50:00 (13 frames). Question 2 of the roadmap frame.
% The primer walks one worked example (bits -> 16-QAM points -> perpendicular axes ->
% IFFT builds the symbol and the grid). No formulas.
% The OTFS / ODDM / delay-Doppler frames that used to follow were removed on 2026-09-13
% (CS audience, no signal-processing background); see parked/03_otfs_oddm_frames.tex.
% The four-frame FFT primer that used to sit between B and C (the 2x2 and 4x4 FFT
% matrices, N^2 vs N log N) was moved on 2026-09-19 to Class_8_WiFi_History/
% 03_first_speed_race.tex, where it follows the 802.11a "64-point FFT" slide. This
% lecture now treats the IFFT/FFT as a labelled box and says so on the C1 frame.

% OFDM primer, 13 frames (2026-09-13, third pass, at the author's request: one idea per
% frame, each motivated by the gap the previous frame leaves; nothing crammed).
%   A. what a radio sends           A1 transition, A2 a wave and its two knobs,
%                                   A3 two knobs = one point, A4 bits pick a point
%   B. many waves, perpendicular    B1 one wave is slow, B2 why not one fast wave (echoes),
%                                   B3 add 64 waves -- how to separate?, B4 perpendicular
%                                   axes, B5 perpendicular waves drawn, B6 the family k/T
%   C. assemble OFDM                C1 transmitter, C2 receiver, C3 the grid

% ---------- A1 ----------
\begin{frame}{Question 2: which signal do we actually use?}
  \vspace{0.3em}
  {\small Question 1 is answered: the two jobs want opposite things, so there is a
   curve of compromises. Whatever signal we pick has to carry \textbf{both} jobs and
   sit somewhere on that curve.\par}

  \vspace{0.45em}
  {\large\textbf{The answer is OFDM} \expand{Orthogonal Frequency Division
   Multiplexing} --- the signal Wi-Fi, 4G and 5G already transmit. No new radio.\par}

  \vspace{0.45em}
  {\small Before we can read OFDM as a radar we need to understand it properly. Most
   of you have only heard the name, so we build it from the ground up, in three parts:\par}
  \vspace{0.2em}
  {\small
  \begin{enumerate}\itemsep0pt
    \item What a radio actually sends --- and why that is a \emph{point}.
    \item Why we send many waves at once, and why they must be \emph{perpendicular}.
    \item Putting it together: transmitter, receiver, and the \emph{grid}.
  \end{enumerate}}

\end{frame}

% ---------- A2 ----------
\begin{frame}{A radio sends a wave, not bits}
  \vspace{0.2em}
  \centering
  \begin{tikzpicture}[scale=0.9,transform shape]
    \draw[-{Latex[length=2mm]},softgray] (0,0) -- (7.6,0) node[anchor=west,font=\scriptsize]{time};
    \draw[very thick,ranblue,domain=0:7.2,samples=160,smooth] plot (\x,{1.0*sin(360*\x/2.4)});
    \draw[very thick,coreorange,domain=0:7.2,samples=160,smooth] plot (\x,{0.55*sin(360*\x/2.4)});
    \draw[very thick,softgreen,dashed,domain=0:7.2,samples=160,smooth] plot (\x,{1.0*sin(360*\x/2.4+90)});
    \draw[<->,softgray] (7.35,0) -- (7.35,1.0) node[midway,right,font=\scriptsize]{how tall};
    \draw[<->,softgray] (0,-1.35) -- (0.6,-1.35) node[midway,below,font=\scriptsize]{where it starts};
  \end{tikzpicture}

  \vspace{0.2em}
  {\scriptsize\textcolor{ranblue}{\textbf{blue}: tall, starts at 0} \quad
   \textcolor{coreorange}{\textbf{orange}: half as tall} \quad
   \textcolor{softgreen}{\textbf{green}: tall, starts a quarter cycle early}\par}

  \vspace{0.6em}
  {\small An antenna cannot emit a 0 or a 1. It emits a \textbf{wave} at some
   frequency. At one frequency a wave has exactly two things you can change:
   \textbf{how tall it is} (amplitude) and \textbf{where in its cycle it starts}
   (phase). Everything else about it is fixed by the frequency.\par}

  \vspace{0.3em}
  \takeaway{Two knobs. To send information, we turn the knobs.}
\end{frame}

% ---------- A3 ----------
\begin{frame}{Two knobs are one point}
  \vspace{0.2em}
  \begin{columns}[T,onlytextwidth]
    \column{0.55\textwidth}
      \centering
      \begin{tikzpicture}[scale=0.78,transform shape]
        \draw[-{Latex[length=2mm]},softgray] (0,0) -- (5.2,0) node[anchor=west,font=\scriptsize]{time};
        \draw[thick,softgray!70,domain=0:4.8,samples=120,smooth] plot (\x,{0.9*cos(360*\x/2.4)});
        \node[font=\scriptsize,text=softgray,anchor=west] at (4.3,0.85) {cosine};
        \draw[thick,softgray!70,dashed,domain=0:4.8,samples=120,smooth] plot (\x,{0.9*sin(360*\x/2.4)});
        \node[font=\scriptsize,text=softgray,anchor=west] at (4.3,-0.75) {sine};
        \draw[very thick,coreorange,domain=0:4.8,samples=120,smooth]
          plot (\x,{-1.35*cos(360*\x/2.4)+1.35*sin(360*\x/2.4)});
        \node[font=\scriptsize\bfseries,text=coreorange,anchor=north] at (2.4,-2.05)
          {orange $= -1.5\times$cosine $+\;1.5\times$sine};
      \end{tikzpicture}
    \column{0.41\textwidth}
      \centering
      \begin{tikzpicture}[scale=0.62,transform shape]
        \draw[-{Latex[length=2mm]},softgray] (-2.6,0) -- (2.6,0) node[anchor=west,font=\scriptsize]{$x$ (cosine)};
        \draw[-{Latex[length=2mm]},softgray] (0,-2.6) -- (0,2.6) node[anchor=south,font=\scriptsize]{$y$ (sine)};
        \draw[dashed,softgray] (-1.5,1.5) -- (-1.5,0) node[anchor=north,font=\scriptsize]{$-1.5$};
        \draw[dashed,softgray] (-1.5,1.5) -- (0,1.5) node[anchor=west,font=\scriptsize]{$1.5$};
        \fill[coreorange] (-1.5,1.5) circle (4.4pt);
        \node[font=\scriptsize\bfseries,text=coreorange,anchor=south east] at (-1.6,1.6) {the wave};
      \end{tikzpicture}
  \end{columns}

  \vspace{0.5em}
  {\small Any wave of the chosen frequency --- any height, any starting point --- can be
   written as ``$x$ parts of a cosine plus $y$ parts of a sine''. So the two knobs are
   just two numbers, and two numbers are a \textbf{point} $(x, y)$.\par}

  \vspace{0.3em}
  \takeaway{One point $=$ one wave. From now on, when you see a dot, think ``a wave''.}
\end{frame}

% ---------- A4 ----------
\begin{frame}{Bits pick the point}
  \vspace{0.1em}
  \begin{columns}[T,onlytextwidth]
    \column{0.52\textwidth}
      {\small Agree on 16 allowed points --- 16 allowed waves (the scheme is called
       16-QAM; the name does not matter). Cut the bits into groups of four:\par}
      \vspace{0.2em}
      {\centering\ttfamily\small 1011\ \ 0010\ \ 1110\ \ 0001\ \ \ldots\par}
      \vspace{0.3em}
      {\small Each group names one point. The radio sends that wave for one
       \textbf{symbol time} $T$, then the next group's wave, and so on.\par}
      \vspace{0.4em}
      {\small The receiver measures $(x, y)$, picks the \textbf{nearest} allowed point,
       and reads four bits back. Noise nudges the measurement, so allowed points must
       keep their distance --- hence 16, not thousands.\par}
    \column{0.44\textwidth}
      \centering
      \begin{tikzpicture}[scale=0.62,transform shape]
        \draw[-{Latex[length=2mm]},softgray] (-2.7,0) -- (2.7,0) node[anchor=west,font=\scriptsize]{$x$};
        \draw[-{Latex[length=2mm]},softgray] (0,-2.7) -- (0,2.7) node[anchor=south,font=\scriptsize]{$y$};
        \foreach \x in {-1.5,-0.5,0.5,1.5}
          \foreach \y in {-1.5,-0.5,0.5,1.5}
            \fill[ranblue!60] (\x,\y) circle (2.6pt);
        \fill[coreorange] (-1.5,1.5) circle (4.4pt);
        \node[font=\scriptsize\bfseries,text=coreorange,anchor=south west] at (-1.38,1.62) {1011};
        \fill[softgray] (-1.3,1.25) circle (1.8pt);
        \node[font=\tiny,text=softgray,anchor=west] at (-1.2,1.15) {measured};
        \node[font=\scriptsize,text=softgray,anchor=north] at (0,-2.85) {16 allowed waves, 4 bits each};
      \end{tikzpicture}
  \end{columns}

  \vspace{0.2em}
  \takeaway{One wave carries 4 bits every $T$: slow.}
\end{frame}

% ---------- B1 ----------
\begin{frame}{Part 2: one wave is slow --- so use many}
  \vspace{0.3em}
  {\small Four bits per symbol time is nowhere near a video stream. Two ways to go
   faster:\par}
  \vspace{0.5em}
  \begin{columns}[T,onlytextwidth]
    \column{0.47\textwidth}
      {\small\textbf{\textcolor{ranpurple}{Option A: one wave, shorter symbols.}}\par}
      \vspace{0.3em}
      {\small Keep one wave and switch points 64 times faster. Symbol time shrinks
       from $T$ to $T/64$.\par}
    \column{0.47\textwidth}
      {\small\textbf{\textcolor{softgreen}{Option B: 64 waves at once.}}\par}
      \vspace{0.3em}
      {\small Use 64 different frequencies at the same time, each carrying its own
       point for the full time $T$. Add them up and send the sum.\par}
  \end{columns}

  \vspace{0.8em}
  {\small Both give 64 times the bits. Every modern radio --- Wi-Fi, 4G, 5G --- chose
   \textbf{B}. The next slide is the reason, and it is the same physics this whole
   lecture is about.\par}

  \vspace{0.3em}
  \takeaway{Same speed. So why does everyone pick many slow waves over one fast one?}
\end{frame}

% ---------- B2 ----------
\begin{frame}{Why not one fast wave? Echoes.}
  \vspace{0.2em}
  \centering
  \begin{tikzpicture}[scale=0.86,transform shape]
    % option A: short symbols and a late echo
    \node[font=\scriptsize\bfseries,text=ranpurple,anchor=east] at (-0.2,1.55) {A: short symbols};
    \foreach \i in {0,...,6}
      \fill[ranpurple!55] (\i*0.9,1.2) rectangle ++(0.8,0.55);
    \foreach \i in {0,...,6}
      \fill[ranpurple!25,draw=ranpurple,dashed] (\i*0.9+0.45,1.2) rectangle ++(0.8,0.55);
    \node[font=\scriptsize,text=ranpurple,anchor=west] at (7.3,1.47) {echo lands on the \emph{next} symbol};
    % option B: one long symbol and the same echo
    \node[font=\scriptsize\bfseries,text=softgreen,anchor=east] at (-0.2,0.0) {B: long symbols};
    \fill[softgreen!55] (0,-0.28) rectangle (6.2,0.27);
    \fill[softgreen!25,draw=softgreen,dashed] (0.45,-0.28) rectangle (6.65,0.27);
    \node[font=\scriptsize,text=softgreen,anchor=west] at (7.3,0.0) {echo mostly lands on \emph{itself}};
    \draw[<->,softgray] (0,-0.75) -- (0.45,-0.75) node[midway,below,font=\scriptsize]{echo delay};
  \end{tikzpicture}

  \vspace{0.6em}
  {\small Radio bounces. A copy of your signal arrives a little late off a wall or a
   building --- about 1\,$\mu$s later for every 300\,m of extra path. With short
   symbols (A) that echo is a smeared copy of the \emph{previous} symbol on top of the
   current one: garbage. With long symbols (B) the echo overlaps the same symbol
   almost entirely; it just makes the wave a bit taller or shorter, which the receiver
   can undo.\par}

  \vspace{0.25em}
  \takeaway{Echoes are why we send many slow waves instead of one fast one.}
\end{frame}

% ---------- B3 ----------
\begin{frame}{Add 64 waves and send the sum --- then what?}
  \vspace{0.2em}
  \centering
  \begin{tikzpicture}[scale=0.88,transform shape]
    \draw[thick,ranblue,domain=0:3.2,samples=120,smooth] plot (\x,{2.6+0.45*sin(360*\x/3.2)});
    \draw[thick,coreorange,domain=0:3.2,samples=120,smooth] plot (\x,{1.7+0.45*sin(360*2*\x/3.2+40)});
    \draw[thick,softgreen,domain=0:3.2,samples=120,smooth] plot (\x,{0.8+0.45*sin(360*3*\x/3.2+110)});
    \node[font=\scriptsize,anchor=west] at (3.3,2.6) {wave 1, its point};
    \node[font=\scriptsize,anchor=west] at (3.3,1.7) {wave 2, its point};
    \node[font=\scriptsize,anchor=west] at (3.3,0.8) {wave 3, its point \ldots\ (64 of them)};
    \node[font=\Large] at (1.6,-0.05) {$+$};
    \draw[very thick,ranpurple,domain=0:3.2,samples=200,smooth]
      plot (\x,{-1.3+0.45*sin(360*\x/3.2)+0.45*sin(360*2*\x/3.2+40)+0.45*sin(360*3*\x/3.2+110)});
    \node[font=\scriptsize\bfseries,text=ranpurple,anchor=west] at (3.3,-1.3) {the sum: this is what the antenna sends};
  \end{tikzpicture}

  \vspace{0.6em}
  {\small The receiver gets only the purple line. It has to recover \textbf{each}
   wave's $(x, y)$ from it --- 64 pairs of numbers out of one squiggle. \\Looking at
   the squiggle, that seems hopeless.\par}

  \vspace{0.25em}
  \takeaway{When can you pull one wave back out of a sum of many?\\That is the question the word ``orthogonal'' answers.}
\end{frame}

% ---------- B4 ----------
\begin{frame}{Perpendicular axes: reading one ignores the other}
  \vspace{0.2em}
  \begin{columns}[T,onlytextwidth]
    \column{0.48\textwidth}
      \centering
      \begin{tikzpicture}[scale=0.66,transform shape]
        \draw[-{Latex[length=2mm]},ranblue,thick] (0,0) -- (3.2,0) node[anchor=west,font=\scriptsize]{$x$};
        \draw[-{Latex[length=2mm]},ranblue,thick] (0,0) -- (0,3.2) node[anchor=south,font=\scriptsize]{$y$};
        \fill[coreorange] (2,2.2) circle (3pt);
        \draw[dashed,softgray] (2,2.2) -- (2,0) node[anchor=north,font=\scriptsize]{2};
        \draw[dashed,softgray] (2,2.2) -- (0,2.2) node[anchor=east,font=\scriptsize]{2.2};
        \node[font=\scriptsize,text=softgreen,align=center,anchor=north] at (1.6,-0.6)
          {perpendicular: reading $x$\\picks up nothing of $y$};
      \end{tikzpicture}
    \column{0.48\textwidth}
      \centering
      \begin{tikzpicture}[scale=0.66,transform shape]
        \draw[-{Latex[length=2mm]},ranpurple,thick] (0,0) -- (3.2,0) node[anchor=west,font=\scriptsize]{$a$};
        \draw[-{Latex[length=2mm]},ranpurple,thick] (0,0) -- (2.77,1.6) node[anchor=west,font=\scriptsize]{$b$};
        \fill[coreorange] (2,2.2) circle (3pt);
        \draw[dashed,softgray] (2,2.2) -- (2,0);
        \draw[dashed,softgray] (2,2.2) -- (2.45,1.42);
        \node[font=\scriptsize,text=ranpurple,align=center,anchor=north] at (1.6,-0.6)
          {tilted: reading $a$\\picks up part of $b$};
      \end{tikzpicture}
  \end{columns}

  \vspace{0.6em}
  {\small Think of a point on a plane. With perpendicular axes, its $x$-reading does
   not change when you slide it along $y$: the two numbers are independent, each can be
   read on its own. Tilt the axes and that breaks --- moving along $b$ changes the
   $a$-reading. The two numbers leak into each other.\par}

  \vspace{0.3em}
  \takeaway{Separable $\Leftrightarrow$ perpendicular. Now we need the same idea for waves.}
\end{frame}

% ---------- B5 ----------
\begin{frame}{Perpendicular waves, drawn}
  \vspace{0.1em}
  \begin{columns}[T,onlytextwidth]
    \column{0.49\textwidth}
      \centering
      {\small\textbf{\textcolor{softgreen}{1 cycle $\times$ 2 cycles in time $T$}}\par}
      \begin{tikzpicture}[scale=0.66,transform shape]
        \draw[softgray] (0,0) -- (4,0);
        \draw[thick,ranblue,domain=0:4,samples=100,smooth] plot (\x,{0.5*sin(360*\x/4)});
        \draw[thick,coreorange,domain=0:4,samples=100,smooth] plot (\x,{0.5*sin(360*2*\x/4)});
        \draw[softgray] (0,-1.9) -- (4,-1.9);
        \fill[softgreen!35,domain=0:4,samples=120,smooth] plot (\x,{-1.9+0.7*sin(360*\x/4)*sin(360*2*\x/4)}) -- (4,-1.9) -- (0,-1.9) -- cycle;
        \draw[thick,softgreen,domain=0:4,samples=120,smooth] plot (\x,{-1.9+0.7*sin(360*\x/4)*sin(360*2*\x/4)});
        \node[font=\scriptsize,anchor=west] at (4.1,-1.9) {product};
        \node[font=\scriptsize\bfseries,text=softgreen,anchor=north] at (2,-2.75) {area above $=$ area below: total 0};
      \end{tikzpicture}
    \column{0.49\textwidth}
      \centering
      {\small\textbf{\textcolor{ranpurple}{1 cycle $\times$ 1.5 cycles in time $T$}}\par}
      \begin{tikzpicture}[scale=0.66,transform shape]
        \draw[softgray] (0,0) -- (4,0);
        \draw[thick,ranblue,domain=0:4,samples=100,smooth] plot (\x,{0.5*sin(360*\x/4)});
        \draw[thick,coreorange,domain=0:4,samples=100,smooth] plot (\x,{0.5*sin(360*1.5*\x/4)});
        \draw[softgray] (0,-1.9) -- (4,-1.9);
        \fill[ranpurple!30,domain=0:4,samples=120,smooth] plot (\x,{-1.9+0.7*sin(360*\x/4)*sin(360*1.5*\x/4)}) -- (4,-1.9) -- (0,-1.9) -- cycle;
        \draw[thick,ranpurple,domain=0:4,samples=120,smooth] plot (\x,{-1.9+0.7*sin(360*\x/4)*sin(360*1.5*\x/4)});
        \node[font=\scriptsize,anchor=west] at (4.1,-1.9) {product};
        \node[font=\scriptsize\bfseries,text=ranpurple,anchor=north] at (2,-2.75) {more above than below: total $\neq 0$};
      \end{tikzpicture}
  \end{columns}

  \vspace{0.5em}
  {\small To measure ``how much of wave 1 is in the sum'', the receiver multiplies the
   sum by wave 1 and adds up over the symbol time $T$. Two waves are
   \textbf{perpendicular} when that product adds up to \textbf{zero}: measuring one
   picks up nothing of the other. Left: it works. Right: it does not --- wave 2 leaks
   into the reading.\par}

  \vspace{0.2em}
  \takeaway{Whole cycles in $T$: perpendicular. Fractional cycles: leak.}
\end{frame}

% ---------- B6 ----------
\begin{frame}{The perpendicular family: frequencies $k/T$}
  \vspace{0.3em}
  \centering
  \begin{tikzpicture}[scale=0.9,transform shape]
    \draw[-{Latex[length=2mm]},softgray] (0,0) -- (9.4,0) node[anchor=west,font=\scriptsize]{frequency};
    \foreach \k/\lab in {1/{$1/T$},2/{$2/T$},3/{$3/T$},4/{$4/T$},5/{$5/T$},6/{$6/T$},7/{$7/T$},8/{$8/T$}}
      { \draw[very thick,ranblue] (\k*1.05,0) -- (\k*1.05,1.4);
        \fill[ranblue] (\k*1.05,1.4) circle (2.4pt);
        \node[font=\scriptsize,anchor=north] at (\k*1.05,-0.08) {\lab}; }
    \node[font=\scriptsize,text=softgray,anchor=west] at (8.7,1.4) {\ldots\ 64 of them};
    \draw[<->,coreorange,thick] (1.05,1.75) -- (2.1,1.75) node[midway,above,font=\scriptsize\bfseries,text=coreorange]{spacing $1/T$};
  \end{tikzpicture}

  \vspace{0.6em}
  {\small A wave fits a whole number of cycles into $T$ exactly when its frequency is
   a multiple of $1/T$. So the waves at $1/T,\ 2/T,\ 3/T,\ldots$ are \textbf{all
   perpendicular to each other}. That is the family OFDM uses. Each member is called
   a \textbf{subcarrier}, and their spacing is $1/T$ --- not a design choice, a
   consequence.\par}

  \vspace{0.3em}
  {\small\textbf{Real numbers.} 5G: $T\approx 67\,\mu$s, so subcarriers 15\,kHz apart
   (30, 60, 120\,kHz variants exist). Wi-Fi: $T = 3.2\,\mu$s, spacing 312.5\,kHz.\par}

  \vspace{0.3em}
  \takeaway{We now know \emph{which} 64 waves to use.}
\end{frame}

% ---------- C1 ----------
% The IFFT/FFT are deliberately left as labelled boxes here: the four frames that
% opened the FFT matrix up (2x2, 4x4, N^2 vs N log N) moved to Class 8 on 2026-09-19.
% Say the sentence below out loud -- a CS audience will accept a named box, but not
% an unnamed one that appears from nowhere.
\begin{frame}{Part 3: the transmitter}
  \vspace{0.2em}
  {\small We now know \emph{which} 64 waves to use. Adding them up --- and, at the far
   end, taking a sum back apart --- is one standard pair of machines: the
   \textbf{IFFT} and the \textbf{FFT}. Today they are labelled boxes; Class 8 opens
   them up and shows you the table inside.\par}

  \vspace{0.45em}
  \centering
  \begin{tikzpicture}[scale=0.84,transform shape,
      pbox/.style={draw=ranblue,fill=blue!5,rounded corners=3pt,minimum width=2.05cm,
        minimum height=0.95cm,align=center,font=\scriptsize}]
    \node[pbox] (bits) at (0,0) {bits\\[-1pt]{\tiny 1011 0010 \ldots}};
    \node[pbox] (qam) at (2.9,0) {64 points\\[-1pt]{\tiny 4 bits each}};
    \node[pbox,draw=coreorange,fill=orange!8] (ifft) at (5.8,0) {IFFT\\[-1pt]{\tiny adds the 64 waves}};
    \node[pbox,minimum width=2.6cm] (sym) at (8.9,0) {one OFDM symbol\\[-1pt]{\tiny the sum, time $T$}};
    \foreach \a/\b in {bits/qam,qam/ifft,ifft/sym} \draw[flow] (\a) -- (\b);
  \end{tikzpicture}

  \vspace{0.7em}
  {\small
  \begin{itemize}\itemsep5pt
    \item Take 256 bits. Cut into 64 groups of four; each group picks a point.
    \item Put point $k$ on the wave at frequency $k/T$ --- on \textbf{subcarrier} $k$.
    \item One IFFT adds the 64 waves into one burst of length $T$: the \textbf{symbol}.
    \item Send it. Then the next 256 bits, the next symbol, and so on.
  \end{itemize}}

  \vspace{0.3em}
  \takeaway{256 bits in, one symbol out --- 64 times what one wave could do.}

  \glossline{FFT = Fast Fourier Transform; IFFT = its inverse. Opened up in Class 8.}
\end{frame}

% ---------- C2 ----------
\begin{frame}{The receiver}
  \vspace{0.3em}
  \centering
  \begin{tikzpicture}[scale=0.84,transform shape,
      pbox/.style={draw=ranblue,fill=blue!5,rounded corners=3pt,minimum width=2.05cm,
        minimum height=0.95cm,align=center,font=\scriptsize}]
    \node[pbox,minimum width=2.6cm] (sym) at (0,0) {received symbol\\[-1pt]{\tiny the sum, plus noise and echoes}};
    \node[pbox,draw=coreorange,fill=orange!8] (fft) at (3.1,0) {FFT\\[-1pt]{\tiny measures each wave}};
    \node[pbox] (qam) at (6.0,0) {64 points\\[-1pt]{\tiny nearest allowed one}};
    \node[pbox] (bits) at (8.9,0) {bits\\[-1pt]{\tiny 1011 0010 \ldots}};
    \foreach \a/\b in {sym/fft,fft/qam,qam/bits} \draw[flow] (\a) -- (\b);
  \end{tikzpicture}

  \vspace{0.7em}
  {\small
  \begin{itemize}\itemsep5pt
    \item One FFT on the received burst: out come 64 $(x, y)$ pairs, one per subcarrier
      --- cleanly, because the waves are perpendicular.
    \item Each pair is nudged by noise; snap it to the nearest allowed point; read 4 bits.
    \item The echo from slide ``Why not one fast wave'' is in here too. It makes each
      subcarrier's point a little larger or rotated --- a per-subcarrier correction,
      then forgotten. \emph{Or} kept: that echo is what the radar half of ISAC reads.
  \end{itemize}}

  \vspace{0.2em}
  \takeaway{Same 64 waves, same FFT, read the other way round.}
\end{frame}

% ---------- C3 ----------
\begin{frame}{Symbol after symbol: the grid}
  \vspace{0.2em}
  \begin{columns}[T,onlytextwidth]
    \column{0.46\textwidth}
      \centering
      \begin{tikzpicture}[scale=0.82,transform shape]
        \foreach \x in {0,...,5}
          \foreach \y in {0,...,3}
            \draw[draw=softgray!50,fill=blue!6] (\x*0.62+0.12,\y*0.52+0.10) rectangle ++(0.58,0.48);
        \fill[orange!45,draw=coreorange] (2*0.62+0.12,2*0.52+0.10) rectangle ++(0.58,0.48);
        \draw[-{Latex[length=2mm]},softgray] (0,0) -- (4.0,0)
          node[anchor=north,font=\scriptsize] {time (one column = one symbol)};
        \draw[-{Latex[length=2mm]},softgray] (0,0) -- (0,2.4);
        \node[font=\scriptsize,rotate=90,anchor=south] at (-0.16,1.2) {frequency (subcarriers)};
      \end{tikzpicture}
    \column{0.50\textwidth}
      \small
      \begin{itemize}\itemsep5pt
        \item Each column is one symbol; each row is one subcarrier.
        \item The orange cell: subcarrier 3 during symbol 3 --- \textbf{one wave, one
          symbol time, one point, four bits}.
        \item 5G calls a cell a \textbf{resource element}. A base station schedules by
          handing out cells: these to you, those to your neighbour.
      \end{itemize}
  \end{columns}

  \vspace{0.4em}
  {\small Every OFDM radio --- Wi-Fi, 4G, 5G --- fills and reads this grid, symbol
   after symbol.\par}

\end{frame}
